\chapter{Aus Welt und Zeit}
\quellseite{7}

Der schreckliche Theaterbrand in Chicago ist in der kirchlichen wie weltlichen Presse viel besprochen worden. Darunter finden sich auch Aussprüche, die das heutige Theater richtig beurtheilen. Ein Blatt, \emph{The Epworth Herald}, sagt nach einer längeren Ausführung: „Wir würden durchaus unsere Pflicht versäumen, wenn wir nicht wieder; wie wir schon oft zuvor gethan haben, die jungen Leute dringend bäten, das Theater zu meiden. Das thun wir nicht wegen der möglichen Gefahr einer Panik oder eines Feuers, sondern wegen des moralischen Schadens, der sicherlich folgt. Theaterleben und Christenleben sind Gegensätze; sie können nicht Hand in Hand gehen. Das eine gewinnt auf Kosten des andern.“ — Bei dem Unglück hat auch ein Prediger der Methodistenkirche sein Leben verloren und ein anderer hat schwere Brandwunden davongetragen. Das Blatt, dem wir dies entnehmen, theilt noch mit: „Es kamen auch viele christliche Leute um, die nur auf Besuch oder in Geschäften in Chicago anwesend waren, die ungekannt und unbekannt glaubten das Schauspielhaus besuchen zu können. In ihren Heimathsgemeinden ahnte niemand, daß sie in ihren sittlichen Ansichten so locker seien.“ Damit wird eine Wahrnehmung ausgesprochen, die auch sonst oft zutrifft. Gar viele, die den Christennamen tragen und sich in ihrer Heimath vor etwas Verwerflichem scheuen würden, thun es umgekehrt, wenn sie in der Fremde sind. Als ob der Allwissende und Allgegenwärtige sie nicht sähe und kennte! — Von dem Schauspiel, das gerade bei dem Brande aufgeführt
\quellseite{8}
wurde, wird gesagt, daß es inhaltlich sonst unanstößig gewesen sei. Mag sein; wir kennen es nicht, wissen aber, daß vieles als unanstößig bezeichnet wird, an dem ein zartes, christliches Gewissen gerechten Anstoß nimmt. Daß aber die Aufführung nicht unanstößig war, zeigt ein Vorfall. Eine der Schauspielerinnen, die eben ihren Theil beendet hatte und von der Bühne abgetreten war, wurde aufgefordert, schnell auf die Straße zu fliehen, um ihr Leben zu retten. „Was? In diesem Kostüm?“ sagte sie. „Ich will fast lieber verbrennen, als mich so auf der Straße zeigen.“ Sie war ohne Zweifel in unzüchtiger Kleidung, wie es auf der Bühne gewöhnlich Brauch ist, vor 2000 Jungen und Alten aufgetreten. Sie schämte sich aber, auf der Straße in solchem Aufputz zu erscheinen. — Wir können nur nochmals an das Schriftwort erinnern, das wie kein anderes auf diesen Vorfall paßt: „Habt nicht lieb die Welt, noch was in der Welt ist. … Die Welt vergehet mit ihrer Lust“, 1 Joh. 2, 15. ff. L. F.

Unter der Zahl der in dem furchtbaren Chicagoer Unglück Umgekommenen, welche nahezu 600 beträgt, befanden sich 39 Lehrerinnen der öffentlichen Schulen und 102 Schulkinder. Diese Zahlen zeigen deutlich, welcher Geist und Sinn unter denen herrscht, welchen man fast allgemein, auch von Seiten der Christen, die Ausbildung und Gestaltung des Charakters der heutigen Jugend anvertraut. Von dem „Geist und Gemüth bildenden und erhebenden“ Einfluß des Theaters erfüllt, ja, davon begeistert, ist es ganz natürlich, daß sie diesen Sinn übertragen auf die ihnen Befohlenen. Und mit welchem Erfolge das geschieht, beweist die Zahl der in jenem Theater gegenwärtig gewesenen Kinder. Sollte ein Christ nicht auch hieraus aufs neue lernen, daß er seine Kinder andern Leuten zur Erziehung und zum Unterricht anzuvertrauen hat als solchen blinden Leitern? Sollen deine Kinder für den Theaterbesuch erzogen werden? Wenn nicht, so sende sie in die christliche Schule. — Einige Zeit nach der entsetzlichen Katastrophe hielt man in Chicago eine Feier ab, in welcher insonderheit der im Theater umgekommenen Lehrer und Schüler gedacht ward. Dabei sagte eine Lehrerin unter anderem in ihrer Rede: „Der Herr rief: ,Lasset die Kindlein zu mir kommen und wehret ihnen nicht, denn solcher ist das Reich Gottes‘ — und sie kamen!“ Das ist eine ganz schändliche Lästerung des heiligen Gotteswortes. Wundern wir uns noch, daß aller Respect vor Gottes Wort fast gänzlich schwindet, ja, auch alles natürliche Verständniß und Gefühl für dasselbe in unserm Volke beinahe völlig erstirbt? Hier hast du eine Erklärung dafür. Andere Erklärungen finden wir in dem greulichen Mißbrauch der Schrift von Seiten so vieler Sectenprediger, desgleichen fast aller Zeitungen, in dem leichtfertigen und spöttischen Anführen von Bibelworten in Gesprächen und dergleichen mehr. Wird es nicht auch hier wieder klar, wo dein Kind nicht hingehört und wo es hingehört? Oder soll dein Kind auch da ankommen, wo die meisten, welche obige Lästerung hörten oder lasen, schon angelangt sind, nämlich bei dem unchristlichen Urtheil: D wie schön!? Jb. S.
